% GENERATED FILE. Edit canonical lessons, then run npm run generate:books.
% canonical-source-hash: fnv1a64:41ba260325c16d8c
% canonical-lessons: MR-C290-havaldar, MR-C290-atak-karne, MR-C290-cauksi-karne, MR-C290-turung, MR-C290-kaidi

\chapter{At the Police Station}
\label{ch:mr-at-the-police-station}
\hlchaptermodality{290}

\begin{chapteropening}
\textbf{By the end of this chapter:} \emph{I can say what the police did: an arrest, questions, a prison.}

Complete the last lesson of chapter 290: I can say what the police did: an arrest, questions, a prison.
\end{chapteropening}

\section[\texorpdfstring{\mr{हवालदार} (havāldār)}{हवालदार (havāldār)}]{\mr{हवालदार} (havāldār) — a constable}

\label{lesson:MR-C290-havaldar}



\begin{quote}
Before the new one: say the Marathi for a pickpocket, then the Marathi for to cheat.
\end{quote}

\subsection*{You'll want to know: \mr{हवालदार}}
\textbf{\mr{हवालदार}} — \emph{havāldār} — \textquotedblleft{}a constable\textquotedblright{}.

\begin{quote}
\mr{आम्हाला} \mr{किती} \mr{वेळ} \mr{थांबावं} \mr{लागेल}, \mr{हे} \mr{हवालदारावर} \mr{अवलंबून} \mr{आहे}.
\end{quote}

\begin{quote}
\emph{āmhālā kitī veḷ thāmbāvaṁ lāgel, he havāldārāvar avlambūn āhe.}
\end{quote}

\begin{quote}
How long we will have to wait depends on the constable.
\end{quote}

Masculine.

\mr{हवालदाराला}, to the constable.

\begin{cousinweb}
\textbf{\mr{हवालदार}} (\emph{havāldār}) is a loan from Persian \emph{ḥavāladār}, 'one entrusted with a charge', usually traced to Arabic \emph{ḥawāla}, 'a charge, a trust'.
\end{cousinweb}

\begin{grammarlens}[title={What you've built}]
One more word for telling what happened.
\end{grammarlens}

\subsection*{Your turn}
\begin{itemize}
\raggedright
  \item \emph{Say it:} \emph{havāldār}
  \item \emph{Say it:} \emph{havāldār}, once more
  \item \emph{Say it:} \emph{āmhālā kitī veḷ thāmbāvaṁ lāgel, he havāldārāvar avlambūn āhe.} — How long we will have to wait depends on the constable.
  \item \emph{Recall:} say the Marathi for a pickpocket, then the Marathi for to cheat, then say \emph{havāldār} again
\end{itemize}

\begin{tcolorbox}[breakable,colback=teal!4,colframe=teal!35!black,title={Before you move on}]
What does \mr{हवालदार} mean? (\textquotedblleft{}A constable\textquotedblright{}.) Say it once more.
\end{tcolorbox}
\section[\texorpdfstring{\mr{अटक} \mr{करणे} (aṭak karṇe)}{अटक करणे (aṭak karṇe)}]{\mr{अटक} \mr{करणे} (aṭak karṇe) — to arrest}

\label{lesson:MR-C290-atak-karne}



\begin{quote}
Before the new one: say the Marathi for to cheat, then the Marathi for a constable.
\end{quote}

\subsection*{You'll want to know: \mr{अटक} \mr{करणे}}
\textbf{\mr{अटक} \mr{करणे}} — \emph{aṭak karṇe} — \textquotedblleft{}to arrest\textquotedblright{}.

\begin{quote}
\mr{पोलिसांनी} \mr{चोराला} \mr{संध्याकाळी} \mr{अटक} \mr{केली}.
\end{quote}

\begin{quote}
\emph{polisānnī corālā sandhyākāḷī aṭak kelī.}
\end{quote}

\begin{quote}
The police arrested the thief in the evening.
\end{quote}

Past, with -\mr{ने} (he …): \mr{त्याने} \mr{अटक} \mr{केली}, \emph{tyāne aṭak kelī}.

The person arrested takes -\mr{ला}; \mr{अटक} is feminine, so the past is \mr{अटक} \mr{केली}.

\begin{cousinweb}
\textbf{\mr{अटक} \mr{करणे}} (\emph{aṭak karṇe}) is built from \emph{aṭak}, 'a hold, a stop', of the family CDIAL lists under \emph{*aṭṭakk}, 'obstruct, stop', with Marathi \emph{aṭakṇe}, 'to be stopped' (CDIAL 182), and \emph{karṇe}, 'to do'.
\end{cousinweb}

\begin{grammarlens}[title={What you've built}]
One more word for telling what happened.
\end{grammarlens}

\subsection*{Your turn}
\begin{itemize}
\raggedright
  \item \emph{Say it:} \emph{aṭak karṇe}
  \item \emph{Say it:} \emph{aṭak karṇe}, once more
  \item \emph{Say it:} \emph{polisānnī corālā sandhyākāḷī aṭak kelī.} — The police arrested the thief in the evening.
  \item \emph{Recall:} say the Marathi for to cheat, then the Marathi for a constable, then say \emph{aṭak karṇe} again
\end{itemize}

\begin{tcolorbox}[breakable,colback=teal!4,colframe=teal!35!black,title={Before you move on}]
What does \mr{अटक} \mr{करणे} mean? (\textquotedblleft{}To arrest\textquotedblright{}.) Say it once more.
\end{tcolorbox}
\section[\texorpdfstring{\mr{चौकशी} \mr{करणे} (caukśī karṇe)}{चौकशी करणे (caukśī karṇe)}]{\mr{चौकशी} \mr{करणे} (caukśī karṇe) — to inquire; to question}

\label{lesson:MR-C290-cauksi-karne}



\begin{quote}
Before the new one: say the Marathi for a constable, then the Marathi for to arrest.
\end{quote}

\subsection*{You'll want to know: \mr{चौकशी} \mr{करणे}}
\textbf{\mr{चौकशी} \mr{करणे}} — \emph{caukśī karṇe} — \textquotedblleft{}to inquire; to question\textquotedblright{}.

\begin{quote}
\mr{हवालदाराने} \mr{दुकानदाराकडे} \mr{चौकशी} \mr{केली}.
\end{quote}

\begin{quote}
\emph{havāldārāne dukāndārākaḍe caukśī kelī.}
\end{quote}

\begin{quote}
The constable made inquiries with the shopkeeper.
\end{quote}

Past, with -\mr{ने} (he …): \mr{त्याने} \mr{चौकशी} \mr{केली}, \emph{tyāne caukśī kelī}.

The person asked takes -\mr{कडे}; the matter takes -\mr{ची}: \mr{अपघाताची} \mr{चौकशी} \mr{करणे}, to look into an accident.

\begin{cousinweb}
\textbf{\mr{चौकशी} \mr{करणे}} (\emph{caukśī karṇe}) joins \emph{caukśī}, 'an inquiry', to \emph{karṇe}, 'to do'; the origin of \emph{caukśī} is not settled.
\end{cousinweb}

\begin{grammarlens}[title={What you've built}]
One more word for telling what happened.
\end{grammarlens}

\subsection*{Your turn}
\begin{itemize}
\raggedright
  \item \emph{Say it:} \emph{caukśī karṇe}
  \item \emph{Say it:} \emph{caukśī karṇe}, once more
  \item \emph{Say it:} \emph{havāldārāne dukāndārākaḍe caukśī kelī.} — The constable made inquiries with the shopkeeper.
  \item \emph{Recall:} say the Marathi for a constable, then the Marathi for to arrest, then say \emph{caukśī karṇe} again
\end{itemize}

\begin{tcolorbox}[breakable,colback=teal!4,colframe=teal!35!black,title={Before you move on}]
What does \mr{चौकशी} \mr{करणे} mean? (\textquotedblleft{}To inquire; to question\textquotedblright{}.) Say it once more.
\end{tcolorbox}
\section[\texorpdfstring{\mr{तुरुंग} (turuṅg)}{तुरुंग (turuṅg)}]{\mr{तुरुंग} (turuṅg) — a prison; a jail}

\label{lesson:MR-C290-turung}



\begin{quote}
Before the new one: say the Marathi for to arrest, then the Marathi for to inquire.
\end{quote}

\subsection*{You'll want to know: \mr{तुरुंग}}
\textbf{\mr{तुरुंग}} — \emph{turuṅg} — \textquotedblleft{}a prison; a jail\textquotedblright{}.

\begin{quote}
\mr{न्यायाधीशांनी} \mr{त्या} \mr{चोराला} \mr{तुरुंगात} \mr{पाठवलं}.
\end{quote}

\begin{quote}
\emph{nyāyādhīśānnī tyā corālā turuṅgāt pāṭhavlaṁ.}
\end{quote}

\begin{quote}
The judge sent that thief to prison.
\end{quote}

Masculine.

\mr{तुरुंगात}, in prison; \mr{तुरुंगातून}, out of prison.

\begin{cousinweb}
\textbf{\mr{तुरुंग}} (\emph{turuṅg}) is the everyday Marathi word for a jail; its origin is not settled.
\end{cousinweb}

\begin{grammarlens}[title={What you've built}]
One more word for telling what happened.
\end{grammarlens}

\subsection*{Your turn}
\begin{itemize}
\raggedright
  \item \emph{Say it:} \emph{turuṅg}
  \item \emph{Say it:} \emph{turuṅg}, once more
  \item \emph{Say it:} \emph{nyāyādhīśānnī tyā corālā turuṅgāt pāṭhavlaṁ.} — The judge sent that thief to prison.
  \item \emph{Recall:} say the Marathi for to arrest, then the Marathi for to inquire, then say \emph{turuṅg} again
\end{itemize}

\begin{tcolorbox}[breakable,colback=teal!4,colframe=teal!35!black,title={Before you move on}]
What does \mr{तुरुंग} mean? (\textquotedblleft{}A prison; a jail\textquotedblright{}.) Say it once more.
\end{tcolorbox}
\section[\texorpdfstring{\mr{कैदी} (kaidī)}{कैदी (kaidī)}]{\mr{कैदी} (kaidī) — a prisoner}

\label{lesson:MR-C290-kaidi}



\begin{quote}
Before the new one: say the Marathi for to inquire, then the Marathi for a prison.
\end{quote}

\subsection*{You'll want to know: \mr{कैदी}}
\textbf{\mr{कैदी}} — \emph{kaidī} — \textquotedblleft{}a prisoner\textquotedblright{}.

\begin{quote}
\mr{एक} \mr{कैदी} \mr{रात्री} \mr{तुरुंगातून} \mr{पळाला}.
\end{quote}

\begin{quote}
\emph{ek kaidī rātrī turuṅgātūn paḷālā.}
\end{quote}

\begin{quote}
A prisoner ran away from the prison at night.
\end{quote}

Masculine.

\mr{कैदी} keeps its form in the plural: \mr{दोन} \mr{कैदी}.

\begin{cousinweb}
\textbf{\mr{कैदी}} (\emph{kaidī}) is a loan from Persian \emph{qaidī}, 'a captive', usually traced to Arabic \emph{qayd}, 'a fetter'.
\end{cousinweb}

\begin{grammarlens}[title={What you've built}]
One more word for telling what happened.
\end{grammarlens}

\subsection*{Your turn}
\begin{itemize}
\raggedright
  \item \emph{Say it:} \emph{kaidī}
  \item \emph{Say it:} \emph{kaidī}, once more
  \item \emph{Say it:} \emph{ek kaidī rātrī turuṅgātūn paḷālā.} — A prisoner ran away from the prison at night.
  \item \emph{Recall:} say the Marathi for to inquire, then the Marathi for a prison, then say \emph{kaidī} again
\end{itemize}

\begin{tcolorbox}[breakable,colback=teal!4,colframe=teal!35!black,title={Before you move on}]
What does \mr{कैदी} mean? (\textquotedblleft{}A prisoner\textquotedblright{}.) Say it once more.
\end{tcolorbox}
